

% -------------------------------------------------
% فصل 2
% -------------------------------------------------

\projectchapter{مبانی نظری و مفاهیم پایه}{
در این فصل مفاهیم نظری و پیش‌زمینه‌های مورد نیاز برای درک روش‌های استفاده‌شده در پروژه بررسی می‌شوند. ابتدا مسئله شناسایی زبان از متن و گفتار معرفی شده و سپس مفاهیم یادگیری نظارت‌شده و بدون نظارت، روش‌های استخراج ویژگی از متن و صوت، الگوریتم‌های طبقه‌بندی و خوشه‌بندی، روش‌های کاهش ابعاد و در نهایت معیارهای ارزیابی مورد استفاده تشریح می‌شوند.
}


% =================================================
\section{مقدمه}
% =================================================

شناسایی خودکار زبان یکی از مسائل مهم در حوزه پردازش زبان طبیعی و پردازش گفتار محسوب می‌شود. هدف اصلی در این مسئله، تعیین زبان یک نمونه ورودی براساس ویژگی‌های قابل استخراج از آن است. ورودی سیستم می‌تواند یک متن یا یک فایل صوتی باشد.

اگر مجموعه زبان‌های موجود به صورت

\begin{equation}
L=\{L_1,L_2,\ldots,L_K\}
\end{equation}

تعریف شود، هدف سیستم یافتن تابعی است که برای یک نمونه ورودی $x$، یکی از زبان‌های مجموعه $L$ را انتخاب کند:

\begin{equation}
f(x)\longrightarrow L_i
\end{equation}

در این پروژه مسئله شناسایی زبان در دو حوزه متن و گفتار بررسی شده است. اگرچه هدف نهایی در هر دو بخش یکسان است، اما نوع داده و نحوه استخراج ویژگی در آن‌ها متفاوت است.

در ادامه این فصل، مفاهیم و روش‌هایی که مستقیماً در پیاده‌سازی پروژه استفاده شده‌اند مورد بررسی قرار می‌گیرند.


% =================================================
\section{شناسایی زبان از متن و گفتار}
% =================================================

\subsection{شناسایی زبان از متن}

در شناسایی زبان مبتنی بر متن\footnote{\lr{Text Language Identification}}، هدف تعیین زبان یک متن براساس ساختار نوشتاری آن است. زبان‌های مختلف از نظر نوع حروف، ساختار واژگان، ترکیب کاراکترها، فراوانی کلمات و ترتیب آن‌ها دارای الگوهای متفاوتی هستند.

حتی زبان‌هایی که از الفبای مشابه استفاده می‌کنند، معمولاً در نحوه قرارگیری حروف و ساختار کلمات تفاوت‌هایی دارند. به همین دلیل می‌توان با استخراج الگوهای آماری از متن، زبان آن را تشخیص داد.

در این پروژه برای نمایش عددی متون از ترکیب ویژگی‌های سطح کلمه و سطح کاراکتر استفاده شده است. این ویژگی‌ها با استفاده از روش
\lr{TF-IDF}
استخراج می‌شوند.


\subsection{شناسایی زبان از گفتار}

در شناسایی زبان از گفتار\footnote{\lr{Spoken Language Identification}}، ورودی سیستم یک سیگنال صوتی است.

یک سیگنال صوتی دیجیتال را می‌توان به صورت دنباله‌ای از نمونه‌ها نمایش داد:

\begin{equation}
x[n], \qquad n=0,1,\ldots,N-1
\end{equation}

که در آن $x[n]$ مقدار دامنه سیگنال در نمونه $n$ است.

استفاده مستقیم از تمام نمونه‌های سیگنال به عنوان ورودی مدل‌های کلاسیک یادگیری ماشین مناسب نیست؛ زیرا هر فایل ممکن است شامل تعداد بسیار زیادی نمونه باشد و همچنین طول فایل‌ها نیز یکسان نباشد.

بنابراین لازم است هر فایل صوتی به یک بردار ویژگی با طول ثابت تبدیل شود. در این پروژه این کار با استخراج مجموعه‌ای از ویژگی‌های طیفی، زمانی و انرژی انجام شده است.


% =================================================
\section{یادگیری ماشین در مسئله شناسایی زبان}
% =================================================

یادگیری ماشین را می‌توان به روش‌های مختلف تقسیم‌بندی کرد. در این پروژه از دو رویکرد اصلی استفاده شده است:

\begin{itemize}
	\item یادگیری نظارت‌شده؛
	\item یادگیری بدون نظارت.
\end{itemize}


% -------------------------------------------------
\subsection{یادگیری نظارت‌شده}
% -------------------------------------------------

در یادگیری نظارت‌شده\footnote{\lr{Supervised Learning}}، برای هر نمونه ورودی یک برچسب صحیح مشخص است.

مجموعه داده آموزشی را می‌توان به صورت زیر نمایش داد:

\begin{equation}
D=
\{(x_1,y_1),(x_2,y_2),\ldots,(x_N,y_N)\}
\end{equation}

که در آن:

\begin{itemize}
	\item $x_i$ بردار ویژگی نمونه $i$؛
	\item $y_i$ برچسب زبان نمونه؛
	\item $N$ تعداد نمونه‌ها است.
\end{itemize}

مدل با استفاده از این داده‌ها رابطه میان ویژگی‌ها و زبان واقعی را یاد گرفته و سپس زبان نمونه‌های جدید را پیش‌بینی می‌کند.

مسئله طبقه‌بندی\footnote{\lr{Classification}} مورد استفاده در این پروژه در این گروه قرار می‌گیرد.


% -------------------------------------------------
\subsection{یادگیری بدون نظارت}
% -------------------------------------------------

در یادگیری بدون نظارت\footnote{\lr{Unsupervised Learning}}، برچسب واقعی نمونه‌ها در اختیار الگوریتم قرار نمی‌گیرد.

در این حالت مجموعه داده به صورت

\begin{equation}
D=\{x_1,x_2,\ldots,x_N\}
\end{equation}

تعریف می‌شود.

الگوریتم تلاش می‌کند نمونه‌های مشابه را براساس ساختار ویژگی‌های آن‌ها در گروه‌هایی قرار دهد. این فرایند در مسئله خوشه‌بندی\footnote{\lr{Clustering}} مورد استفاده قرار می‌گیرد.

در این پروژه، خوشه‌بندی برای بررسی این موضوع انجام شده است که آیا نمونه‌های متعلق به یک زبان، بدون در اختیار داشتن برچسب زبان، به صورت طبیعی در کنار یکدیگر قرار می‌گیرند یا خیر.


% =================================================
\section{استخراج ویژگی از داده‌های متنی}
% =================================================

الگوریتم‌های کلاسیک یادگیری ماشین نمی‌توانند متن خام را مستقیماً پردازش کنند. بنابراین لازم است متن به یک نمایش عددی تبدیل شود.

در این پروژه از روش
\lr{TF-IDF}
در کنار
\lr{Word N-Gram}
و
\lr{Character N-Gram}
استفاده شده است.


% -------------------------------------------------
\subsection{مفهوم \lr{N-Gram}}
% -------------------------------------------------

\lr{N-Gram}
به دنباله‌ای شامل $n$ عنصر متوالی گفته می‌شود. این عناصر می‌توانند کلمه یا کاراکتر باشند.

برای مثال، در کلمه

\begin{center}
\lr{language}
\end{center}

برخی
\lr{Character Bigram}
ها عبارت‌اند از:

\begin{center}
\lr{la,\ an,\ ng,\ gu,\ ua,\ ag,\ ge}
\end{center}

همچنین برای جمله

\begin{center}
\lr{this is a language}
\end{center}

برخی
\lr{Word Bigram}
ها شامل موارد زیر هستند:

\begin{center}
\lr{this is,\ is a,\ a language}
\end{center}

در پروژه حاضر از ویژگی‌های زیر استفاده شده است:

\begin{itemize}
	\item \lr{Character N-Gram} با مرتبه‌های 2 تا 4؛
	\item \lr{Word N-Gram} با مرتبه‌های 1 تا 2.
\end{itemize}

استفاده هم‌زمان از این دو نوع ویژگی باعث می‌شود هم اطلاعات مربوط به ساختار کلمات و هم الگوهای موجود در ترکیب حروف در اختیار مدل قرار گیرد.


% -------------------------------------------------
\subsection{روش \lr{TF-IDF}}
% -------------------------------------------------

روش
\lr{TF-IDF}\footnote{\lr{Term Frequency--Inverse Document Frequency}}
یکی از روش‌های متداول برای تبدیل متن به نمایش عددی است.

این روش از دو مؤلفه اصلی تشکیل شده است.


\subsubsection{\lr{Term Frequency}}

مؤلفه
\lr{Term Frequency}
میزان تکرار یک ویژگی $t$ در سند $d$ را اندازه‌گیری می‌کند.

در حالت ساده می‌توان نوشت:

\begin{equation}
TF(t,d)=\mathrm{count}(t,d)
\end{equation}

در پروژه حاضر از
\lr{Sublinear TF}
استفاده شده است. در این حالت برای فراوانی‌های غیرصفر داریم:

\begin{equation}
TF(t,d)=1+\log(\mathrm{count}(t,d))
\end{equation}

استفاده از لگاریتم باعث می‌شود تکرار بسیار زیاد یک ویژگی، اثر بیش از حدی بر مدل ایجاد نکند.


\subsubsection{\lr{Inverse Document Frequency}}

مؤلفه
\lr{IDF}
نشان می‌دهد یک ویژگی در مجموعه اسناد تا چه اندازه خاص یا عمومی است.

فرم هموارشده آن را می‌توان به صورت زیر نمایش داد:

\begin{equation}
IDF(t)
=
\log
\left(
\frac{1+N}{1+df(t)}
\right)
+1
\end{equation}

که در آن:

\begin{itemize}
	\item $N$ تعداد کل اسناد؛
	\item $df(t)$ تعداد اسنادی است که ویژگی $t$ در آن‌ها وجود دارد.
\end{itemize}

در نهایت مقدار
\lr{TF-IDF}
از رابطه زیر به دست می‌آید:

\begin{equation}
TFIDF(t,d)
=
TF(t,d)
\times
IDF(t)
\end{equation}

بنابراین ویژگی‌هایی که در یک متن اهمیت زیادی دارند ولی در تمام متون به شکل گسترده تکرار نمی‌شوند، وزن بیشتری دریافت می‌کنند.


% -------------------------------------------------
\subsection{ترکیب ویژگی‌های کاراکتری و کلمه‌ای}
% -------------------------------------------------

در پروژه، ویژگی‌های استخراج‌شده در سطح کاراکتر و کلمه با یکدیگر ترکیب شده‌اند.

اگر ماتریس ویژگی‌های کاراکتری را با $X_c$ و ماتریس ویژگی‌های واژگانی را با $X_w$ نمایش دهیم، نمایش ترکیبی به‌صورت زیر خواهد بود:

\begin{equation}
X=
\left[
X_c
\quad
X_w
\right]
\end{equation}

این ترکیب باعث می‌شود اطلاعات مربوط به ساختار کلمات و الگوهای موجود در توالی کاراکترها به‌صورت هم‌زمان در اختیار مدل قرار گیرد. تنظیمات دقیق بردارسازهای کاراکتری و واژگانی در فصل سوم ارائه خواهند شد.


% -------------------------------------------------
\subsection{انتخاب ویژگی با \lr{Chi-Square}}
% -------------------------------------------------

برای انتخاب ویژگی‌های مناسب در بخش طبقه‌بندی متن از آزمون
\lr{Chi-Square}
استفاده شده است.

رابطه کلی آزمون کای‌دو به‌صورت زیر تعریف می‌شود:

\begin{equation}
\chi^2
=
\sum
\frac{(O-E)^2}{E}
\end{equation}

که در آن:

\begin{itemize}

	\item $O$ مقدار مشاهده‌شده؛

	\item $E$ مقدار مورد انتظار است.

\end{itemize}

هرچه مقدار
$\chi^2$
برای یک ویژگی بیشتر باشد، وابستگی آماری آن ویژگی به برچسب زبان بیشتر است.

ازآنجاکه آزمون
\lr{Chi-Square}
به برچسب زبان نیاز دارد، این مرحله فقط در مسیر طبقه‌بندی متن استفاده شده و در خوشه‌بندی متن به کار نرفته است.

برای جلوگیری از نشت اطلاعات، انتخاب ویژگی باید فقط با استفاده از بخش آموزشی انجام شود. تعداد دقیق ویژگی‌های انتخاب‌شده و نحوه اجرای این مرحله در فصل سوم تشریح خواهند شد.

% =================================================
\section{استخراج ویژگی از داده‌های صوتی}
% =================================================

هدف از استخراج ویژگی صوتی، تبدیل سیگنال خام با طول متغیر به یک بردار عددی با طول ثابت است.

خانواده ویژگی‌های صوتی مورد استفاده در ادامه معرفی می‌شوند و تعداد و ساختار دقیق ویژگی‌های نسخه نهایی در فصل سوم ارائه خواهد شد.


% -------------------------------------------------
\subsection{ضرایب \lr{MFCC}}
% -------------------------------------------------

یکی از مهم‌ترین ویژگی‌های مورد استفاده در پردازش گفتار،
\lr{MFCC}\footnote{\lr{Mel-Frequency Cepstral Coefficients}}
است.

مقیاس
\lr{Mel}
برای نزدیک کردن نمایش فرکانس به نحوه درک انسان از صوت مورد استفاده قرار می‌گیرد.

تبدیل تقریبی فرکانس معمولی به مقیاس
\lr{Mel}
به صورت زیر است:

\begin{equation}
Mel(f)
=
2595
\log_{10}
\left(
1+\frac{f}{700}
\right)
\end{equation}

پس از محاسبه طیف و عبور آن از بانک فیلترهای
\lr{Mel}،
از لگاریتم انرژی فیلترها و سپس تبدیل کسینوسی گسسته استفاده می‌شود.

یک فرم کلی برای محاسبه ضریب
\lr{MFCC}
به صورت زیر است:

\begin{equation}
c_n
=
\sum_{k=1}^{K}
\log(M_k)
\cos
\left[
\frac{\pi n}{K}
\left(
k-\frac{1}{2}
\right)
\right]
\end{equation}

که در آن:

\begin{itemize}
	\item $M_k$ انرژی فیلتر \lr{Mel} شماره $k$؛
	\item $K$ تعداد فیلترها؛
	\item $c_n$ ضریب \lr{MFCC} است.
\end{itemize}

از آنجا که
\lr{MFCC}
برای فریم‌های مختلف سیگنال محاسبه می‌شود، برای تبدیل آن به ویژگی با طول ثابت، میانگین و انحراف معیار هر ضریب استخراج شده است.

میانگین ضریب $i$ برابر است با:

\begin{equation}
\mu_i
=
\frac{1}{T}
\sum_{t=1}^{T}
c_i(t)
\end{equation}

و انحراف معیار آن برابر است با:

\begin{equation}
\sigma_i
=
\sqrt{
\frac{1}{T}
\sum_{t=1}^{T}
(c_i(t)-\mu_i)^2
}
\end{equation}

% -------------------------------------------------
\subsection{\lr{Delta MFCC} و \lr{Delta-Delta MFCC}}
% -------------------------------------------------

\lr{MFCC}
وضعیت طیفی سیگنال را در هر لحظه نمایش می‌دهد، اما نحوه تغییر این ضرایب در طول زمان نیز می‌تواند اطلاعات مهمی در مورد گفتار ارائه دهد.

برای ثبت تغییرات زمانی از ویژگی‌های
\lr{Delta MFCC}
و
\lr{Delta-Delta MFCC}
استفاده شده است.

یک تقریب متداول برای
\lr{Delta}
به صورت زیر است:

\begin{equation}
\Delta c_t
=
\frac{
\sum_{n=1}^{N}
n(c_{t+n}-c_{t-n})
}{
2
\sum_{n=1}^{N}
n^2
}
\end{equation}

\lr{Delta}
بیانگر سرعت تغییر
\lr{MFCC}
و
\lr{Delta-Delta}
بیانگر تغییرات مرتبه دوم یا شتاب تغییر ضرایب است.

برای ضرایب
\lr{Delta}
و
\lr{Delta-Delta}
نیز میانگین و انحراف معیار محاسبه شده است.


% -------------------------------------------------
\subsection{\lr{Zero Crossing Rate}}
% -------------------------------------------------

نرخ عبور از صفر یا
\lr{Zero Crossing Rate}
تعداد دفعاتی را اندازه‌گیری می‌کند که علامت سیگنال در طول زمان تغییر می‌کند.

رابطه آن را می‌توان به صورت زیر نمایش داد:

\begin{equation}
ZCR
=
\frac{1}{2(N-1)}
\sum_{n=1}^{N-1}
\left|
\mathrm{sgn}(x[n])
-
\mathrm{sgn}(x[n-1])
\right|
\end{equation}

برای ساخت یک نمایش با طول ثابت می‌توان از میانگین و انحراف معیار
\lr{ZCR}
در فریم‌های مختلف استفاده کرد.


% -------------------------------------------------
\subsection{ویژگی‌های طیفی}
% -------------------------------------------------

علاوه بر
\lr{MFCC}،
تعدادی ویژگی برای توصیف شکل و توزیع انرژی طیف استخراج شده‌اند.


\subsubsection{\lr{Spectral Centroid}}

\lr{Spectral Centroid}
مرکز جرم طیف فرکانسی را نشان می‌دهد:

\begin{equation}
SC
=
\frac{
\sum_k f_k S_k
}{
\sum_k S_k
}
\end{equation}

که $f_k$ فرکانس و $S_k$ مقدار طیفی متناظر با آن است.

مقادیر بالاتر این ویژگی به معنی تمرکز بیشتر انرژی در فرکانس‌های بالاتر است.


\subsubsection{\lr{Spectral Rolloff}}

\lr{Spectral Rolloff}
فرکانسی است که درصد مشخصی از انرژی کل طیف در زیر آن قرار دارد.

اگر $f_r$ فرکانس
\lr{Rolloff}
باشد:

\begin{equation}
\sum_{f_k\leq f_r}S_k
\geq
\alpha
\sum_k S_k
\end{equation}

که $\alpha$ نسبت مشخصی از انرژی کل طیف است.


\subsubsection{\lr{Spectral Bandwidth}}

\lr{Spectral Bandwidth}
میزان پراکندگی طیف در اطراف
\lr{Spectral Centroid}
را مشخص می‌کند:

\begin{equation}
BW
=
\sqrt{
\frac{
\sum_k
S_k(f_k-SC)^2
}{
\sum_kS_k
}
}
\end{equation}


\subsubsection{\lr{Spectral Flatness}}

\lr{Spectral Flatness}
میزان تخت یا قله‌دار بودن طیف را اندازه‌گیری می‌کند.

این معیار از نسبت میانگین هندسی به میانگین حسابی طیف به دست می‌آید:

\begin{equation}
SF
=
\frac{
\left(
\prod_{k=1}^{K}S_k
\right)^{1/K}
}{
\frac{1}{K}
\sum_{k=1}^{K}S_k
}
\end{equation}

این ویژگی می‌تواند اطلاعاتی درباره میزان نویزگونه یا ساختاریافته بودن طیف ارائه دهد.


% -------------------------------------------------
\subsection{\lr{RMS Energy}}
% -------------------------------------------------

انرژی
\lr{RMS}
یا
\lr{Root Mean Square Energy}
برای اندازه‌گیری سطح انرژی سیگنال استفاده می‌شود.

رابطه آن به صورت زیر است:

\begin{equation}
RMS
=
\sqrt{
\frac{1}{N}
\sum_{n=0}^{N-1}
x[n]^2
}
\end{equation}

برای ساخت ویژگی‌های سطح فایل می‌توان از میانگین و انحراف معیار
\lr{RMS}
استفاده کرد.


% -------------------------------------------------
\subsection{مدت زمان فایل صوتی}
% -------------------------------------------------

مدت زمان فایل صوتی نیز به عنوان یک ویژگی در نظر گرفته شده است.

اگر تعداد نمونه‌های سیگنال برابر $N$ و نرخ نمونه‌برداری برابر $f_s$ باشد، مدت زمان از رابطه زیر به دست می‌آید:

\begin{equation}
Duration
=
\frac{N}{f_s}
\end{equation}


% =================================================
\section{الگوریتم‌های طبقه‌بندی}
% =================================================

در مسیر متن، الگوریتم‌های
\lr{K-Nearest Neighbors}،
\lr{Decision Tree}،
\lr{Logistic Regression}
و
\lr{Multinomial Naive Bayes}
بررسی شده‌اند. در مسیر گفتار نیز الگوریتم‌های
\lr{Logistic Regression}،
\lr{Support Vector Machine}،
\lr{Random Forest}
و
\lr{K-Nearest Neighbors}
مورد استفاده قرار گرفته‌اند. جزئیات تنظیمات و ابرپارامترهای این مدل‌ها در فصل چهارم ارائه خواهند شد.


% -------------------------------------------------
\subsection{\lr{Logistic Regression}}
% -------------------------------------------------

رگرسیون لجستیک\footnote{\lr{Logistic Regression}} با وجود نام آن، یکی از الگوریتم‌های متداول برای طبقه‌بندی است.

در مسئله چندکلاسه، احتمال تعلق نمونه $x$ به کلاس $k$ را می‌توان با تابع
\lr{Softmax}
نمایش داد:

\begin{equation}
P(y=k|x)
=
\frac{
e^{w_k^Tx+b_k}
}{
\sum_{j=1}^{K}
e^{w_j^Tx+b_j}
}
\end{equation}

کلاس دارای بیشترین احتمال به عنوان خروجی مدل انتخاب می‌شود:

\begin{equation}
\hat{y}
=
\arg\max_k
P(y=k|x)
\end{equation}

در هر دو مسیر متن و گفتار،
\lr{Logistic Regression}
به‌عنوان یک طبقه‌بند خطی مورد بررسی قرار گرفته است. منظم‌سازی برای کنترل پیچیدگی مدل و کاهش بیش‌برازش به کار می‌رود.


% -------------------------------------------------
\subsection{\lr{Multinomial Naive Bayes}}
% -------------------------------------------------

روش
\lr{Naive Bayes}
یک الگوریتم احتمالاتی مبتنی بر قضیه بیز است:

\begin{equation}
P(C|X)
=
\frac{
P(X|C)P(C)
}{
P(X)
}
\end{equation}

مدل
\lr{Multinomial Naive Bayes}
به‌ویژه برای داده‌های متنی و ویژگی‌های مبتنی بر فراوانی یا
\lr{TF-IDF}
کاربرد دارد.

در این روش از پارامتر هموارسازی
$\alpha$
برای جلوگیری از صفر شدن احتمال برخی ویژگی‌ها استفاده می‌شود.

این الگوریتم در پروژه برای طبقه‌بندی داده‌های متنی استفاده شده است.


% -------------------------------------------------
\subsection{\lr{Decision Tree}}
% -------------------------------------------------

درخت تصمیم\footnote{\lr{Decision Tree}} با تقسیم پیاپی فضای ویژگی‌ها، نمونه‌ها را به گروه‌های خالص‌تر تقسیم می‌کند.

یکی از معیارهای رایج برای محاسبه ناخالصی یک گره،
\lr{Gini Impurity}
است:

\begin{equation}
Gini
=
1-
\sum_{k=1}^{K}
p_k^2
\end{equation}

که $p_k$ نسبت نمونه‌های کلاس $k$ در گره مورد نظر است.

این الگوریتم در پروژه برای طبقه‌بندی متن استفاده شده است. کنترل پیچیدگی درخت می‌تواند احتمال بیش‌برازش را کاهش دهد.


% -------------------------------------------------
\subsection{\lr{Support Vector Machine}}
% -------------------------------------------------

ماشین بردار پشتیبان یا
\lr{SVM}\footnote{\lr{Support Vector Machine}}
یکی از الگوریتم‌های اصلی مورد استفاده در بخش گفتار است.

در حالت خطی، تابع تصمیم به صورت زیر تعریف می‌شود:

\begin{equation}
f(x)
=
w^\top x+b
\end{equation}

هدف
\lr{SVM}
یافتن مرزی است که فاصله میان نزدیک‌ترین نمونه‌های کلاس‌ها و مرز تصمیم را بیشینه کند.

برای داده‌هایی که کاملاً جداشدنی نیستند، از مفهوم حاشیه نرم استفاده می‌شود. تابع هدف آن را می‌توان به‌صورت زیر نوشت:

\begin{equation}
\min_{w,b,\xi}
\left(
\frac{1}{2}\|w\|_2^2
+
C\sum_{i=1}^{N}\xi_i
\right)
\quad
\text{s.t.}
\quad
y_i(w^\top x_i+b)\geq 1-\xi_i,
\quad
\xi_i\geq 0
\end{equation}

پارامتر $C$ تعادل میان بزرگ‌بودن حاشیه و جریمه خطاهای آموزشی را کنترل می‌کند.

برای جداسازی غیرخطی می‌توان از تابع کرنل استفاده کرد. کرنل مورد استفاده در این پروژه،
\lr{RBF}
است:

\begin{equation}
K(x_i,x_j)
=
\exp
\left(
-\gamma
\|x_i-x_j\|_2^2
\right)
\end{equation}

پارامتر $\gamma$ محدوده اثر هر نمونه را کنترل می‌کند. مقادیر بررسی‌شده برای $C$ و $\gamma$ در فصل چهارم ارائه خواهند شد.


% -------------------------------------------------
\subsection{\lr{K-Nearest Neighbors}}
% -------------------------------------------------

الگوریتم
\lr{KNN}\footnote{\lr{K-Nearest Neighbors}}
کلاس نمونه جدید را براساس نزدیک‌ترین نمونه‌های آموزشی مشخص می‌کند.

یکی از معیارهای فاصله، فاصله اقلیدسی است:

\begin{equation}
d(x_i,x_j)
=
\sqrt{
\sum_{m=1}^{p}
(x_{im}-x_{jm})^2
}
\end{equation}

پس از محاسبه فاصله، $K$ نزدیک‌ترین همسایه انتخاب شده و کلاس براساس رأی آن‌ها تعیین می‌شود.

برای نمایش‌های متنی می‌توان از فاصله کسینوسی استفاده کرد:

\begin{equation}
d_{\mathrm{cosine}}(x_i,x_j)
=
1-
\frac{x_i^\top x_j}{\|x_i\|_2\|x_j\|_2}
\end{equation}

معیار فاصله و نحوه وزن‌دهی همسایه‌ها باید متناسب با نوع نمایش داده انتخاب شوند. جزئیات انتخاب‌های انجام‌شده در دو مسیر در فصل چهارم بیان خواهند شد.


% -------------------------------------------------
\subsection{\lr{Random Forest}}
% -------------------------------------------------

جنگل تصادفی\footnote{\lr{Random Forest}} یک روش
\lr{Ensemble}
است که از مجموعه‌ای از درخت‌های تصمیم استفاده می‌کند.

خروجی نهایی در مسئله طبقه‌بندی از رأی اکثریت درخت‌ها به دست می‌آید:

\begin{equation}
\hat{y}
=
\mathrm{mode}
\{
h_1(x),
h_2(x),
\ldots,
h_M(x)
\}
\end{equation}

استفاده از چندین درخت باعث کاهش حساسیت مدل به نویز و کاهش احتمال بیش‌برازش نسبت به یک درخت منفرد می‌شود.

این الگوریتم در پروژه برای طبقه‌بندی ویژگی‌های گفتاری استفاده شده است.


% =================================================
\section{الگوریتم‌های خوشه‌بندی}
% =================================================

در فرآیند خوشه‌بندی، برچسب واقعی زبان‌ها در مرحله آموزش الگوریتم‌ها استفاده نشده و این برچسب‌ها تنها پس از تشکیل خوشه‌ها برای ارزیابی کیفیت تفکیک زبان‌ها به کار گرفته شده‌اند.

% -------------------------------------------------
\subsection{\lr{K-Means}}
% -------------------------------------------------

الگوریتم
\lr{K-Means}
داده‌ها را به $K$ خوشه تقسیم می‌کند و تلاش می‌کند فاصله نمونه‌ها از مرکز خوشه مربوط به خود را حداقل کند.

تابع هدف به صورت زیر است:

\begin{equation}
J
=
\sum_{k=1}^{K}
\sum_{x_i\in C_k}
\|x_i-\mu_k\|^2
\end{equation}

که در آن:

\begin{itemize}
	\item $C_k$ خوشه شماره $k$؛
	\item $\mu_k$ مرکز خوشه است.
\end{itemize}


% -------------------------------------------------
\subsection{\lr{Hierarchical Clustering}}
% -------------------------------------------------

در
\lr{Agglomerative Hierarchical Clustering}
ابتدا هر نمونه به عنوان یک خوشه مستقل در نظر گرفته می‌شود.

سپس در هر مرحله دو خوشه نزدیک‌تر با یکدیگر ادغام می‌شوند تا در نهایت تعداد مورد نظر خوشه‌ها حاصل شود.

در روش
\lr{Ward Linkage}
ادغامی انتخاب می‌شود که کمترین افزایش را در واریانس داخل خوشه‌ها ایجاد کند.


% -------------------------------------------------
\subsection{\lr{Gaussian Mixture Model}}
% -------------------------------------------------

در بخش متنی و گفتاری پروژه از
\lr{Gaussian Mixture Model}
نیز استفاده شده است.

در این روش فرض می‌شود توزیع داده از ترکیب چند توزیع
\lr{Gaussian}
تشکیل شده باشد:

\begin{equation}
p(x)
=
\sum_{k=1}^{K}
\pi_k
\mathcal{N}
(x|\mu_k,\Sigma_k)
\end{equation}

که در آن:

\begin{itemize}
	\item $\pi_k$ وزن مؤلفه $k$؛
	\item $\mu_k$ بردار میانگین؛
	\item $\Sigma_k$ ماتریس کوواریانس است.
\end{itemize}

این روش برخلاف
\lr{K-Means}
امکان محاسبه احتمال تعلق هر نمونه به خوشه‌های مختلف را فراهم می‌کند.

یکی از معیارهای قابل استفاده برای انتخاب پیچیدگی مدل مخلوط گاوسی، معیار اطلاعات بیزی یا
\lr{BIC}\footnote{\lr{Bayesian Information Criterion}}
است:

\begin{equation}
BIC
=
-2\log(\hat{\mathcal{L}})
+
q\log(N)
\end{equation}

که در آن $\hat{\mathcal{L}}$ بیشینه درست‌نمایی و $q$ تعداد پارامترهای مدل است. مقدار کمتر این معیار نشان‌دهنده تعادل مناسب‌تر میان برازش و پیچیدگی مدل است.


% -------------------------------------------------
\subsection{\lr{DBSCAN}}
% -------------------------------------------------

الگوریتم
\lr{DBSCAN}\footnote{\lr{Density-Based Spatial Clustering of Applications with Noise}}
یک روش خوشه‌بندی مبتنی بر چگالی است.

دو پارامتر اصلی آن عبارت‌اند از:

\begin{itemize}
	\item \lr{eps} یا شعاع همسایگی؛
	\item \lr{min\_samples} یا حداقل تعداد نقاط همسایه.
\end{itemize}

یک نقطه در صورتی نقطه مرکزی محسوب می‌شود که:

\begin{equation}
|N_{\varepsilon}(x)|
\geq
MinPts
\end{equation}

یکی از مزایای مهم
\lr{DBSCAN}
قابلیت تشخیص نمونه‌های نویزی است که معمولاً با برچسب $-1$ مشخص می‌شوند.

% =================================================
\section{آماده‌سازی و کاهش ابعاد ویژگی‌ها}
% =================================================


% -------------------------------------------------
\subsection{استانداردسازی ویژگی‌ها}
% -------------------------------------------------

ویژگی‌های صوتی دارای مقیاس‌های عددی متفاوت هستند. این موضوع می‌تواند به‌ویژه بر الگوریتم‌های مبتنی بر فاصله اثرگذار باشد.

استانداردسازی یک ویژگی به‌صورت زیر تعریف می‌شود:

\begin{equation}
z
=
\frac{x-\mu}{\sigma}
\end{equation}

که در آن:

\begin{itemize}
	\item $\mu$ میانگین ویژگی؛
	\item $\sigma$ انحراف معیار ویژگی است.
\end{itemize}

پس از استانداردسازی، ویژگی‌ها دارای میانگین تقریباً صفر و انحراف معیار تقریباً یک خواهند بود.


% -------------------------------------------------
\subsection{\lr{Principal Component Analysis}}
% -------------------------------------------------

روش
\lr{PCA}\footnote{\lr{Principal Component Analysis}}
برای کاهش ابعاد و حذف هم‌بستگی میان ویژگی‌ها استفاده می‌شود.

اگر ماتریس کوواریانس داده را با $\Sigma$ نمایش دهیم:

\begin{equation}
\Sigma v_i
=
\lambda_i v_i
\end{equation}

که در آن:

\begin{itemize}
	\item $v_i$ بردار ویژه؛
	\item $\lambda_i$ مقدار ویژه متناظر است.
\end{itemize}

مؤلفه‌هایی که مقدار ویژه بزرگ‌تری دارند، بخش بیشتری از واریانس داده را توضیح می‌دهند.

در پروژه حاضر، این روش برای ایجاد نمایش کم‌بعد ویژگی‌های گفتاری در مسیر خوشه‌بندی استفاده شده است. تعداد مؤلفه‌ها و جزئیات اجرا در فصل سوم ارائه خواهند شد.


% -------------------------------------------------
\subsection{\lr{Truncated SVD}}
% -------------------------------------------------

ماتریس
\lr{TF-IDF}
معمولاً دارای تعداد ابعاد بالا و ساختاری
\lr{Sparse}
است.

در خوشه‌بندی متن از
\lr{Truncated SVD}
برای کاهش ابعاد فضای ویژگی استفاده شده است.

تجزیه
\lr{SVD}
به صورت زیر تعریف می‌شود:

\begin{equation}
X
=
U\Sigma V^T
\end{equation}

در نسخه
\lr{Truncated}
تنها تعداد محدودی از مؤلفه‌ها نگه داشته می‌شوند:

\begin{equation}
X
\approx
U_r
\Sigma_r
V_r^T
\end{equation}

این روش امکان کاهش ابعاد ماتریس
\lr{TF-IDF}
را بدون تبدیل مستقیم آن به یک ماتریس
\lr{Dense}
فراهم می‌کند.

در پروژه حاضر، این روش برای ایجاد نمایش کم‌بعد متن در مسیر خوشه‌بندی استفاده شده است. تعداد مؤلفه‌ها و نحوه نرمال‌سازی نمایش نهایی در فصل سوم بیان خواهند شد.


% =================================================
\section{تنظیم و اعتبارسنجی مدل}
% =================================================


% -------------------------------------------------
\subsection{\lr{Cross Validation}}
% -------------------------------------------------

برای ارزیابی قابل اعتمادتر مدل، از
\lr{Cross Validation}
استفاده شده است.

در این پروژه برای جلوگیری از قرار گرفتن نمونه‌های وابسته به یک منبع مشترک در بخش‌های مختلف آموزش و ارزیابی، از روش
\lr{Stratified Group K-Fold}
استفاده شده است.

در این روش علاوه بر حفظ نسبت نمونه‌های کلاس‌های مختلف در
\lr{Fold}
ها، نمونه‌های متعلق به یک گروه مشترک در یک بخش قرار می‌گیرند.

اگر امتیاز هر
\lr{Fold}
برابر $S_i$ باشد، امتیاز میانگین به صورت زیر محاسبه می‌شود:

\begin{equation}
CV
=
\frac{1}{K}
\sum_{i=1}^{K}
S_i
\end{equation}


% -------------------------------------------------
\subsection{\lr{Grid Search}}
% -------------------------------------------------

برای انتخاب بهترین مقادیر ابرپارامترها می‌توان از روش
\lr{Grid Search}
استفاده کرد.

این روش ترکیب‌های مختلف پارامترهای تعیین‌شده را آزمایش کرده و هر ترکیب را با استفاده از
\lr{Cross Validation}
ارزیابی می‌کند.

در این پروژه، معیارهایی مانند
\lr{Macro F1}
برای مقایسه مدل‌ها و انتخاب تنظیمات مناسب مورد استفاده قرار گرفته‌اند.

تعداد بخش‌های اعتبارسنجی و شبکه ابرپارامترهای هر مدل در فصل چهارم ارائه خواهند شد.

% =================================================
\section{معیارهای ارزیابی طبقه‌بندی}
% =================================================


% -------------------------------------------------
\subsection{\lr{Accuracy}}
% -------------------------------------------------

دقت کلی یا
\lr{Accuracy}
نسبت تعداد پیش‌بینی‌های صحیح به کل پیش‌بینی‌ها است. در حالت چندکلاسه می‌توان نوشت:

\begin{equation}
Accuracy
=
\frac{1}{N}
\sum_{i=1}^{N}
\mathbb{I}(y_i=\hat{y}_i)
\end{equation}


% -------------------------------------------------
\subsection{\lr{Precision}}
% -------------------------------------------------

معیار
\lr{Precision}
مشخص می‌کند از میان نمونه‌هایی که مدل به یک کلاس نسبت داده است، چه نسبتی واقعاً متعلق به آن کلاس بوده‌اند:

\begin{equation}
Precision
=
\frac{TP}{TP+FP}
\end{equation}


% -------------------------------------------------
\subsection{\lr{Recall}}
% -------------------------------------------------

\lr{Recall}
نشان می‌دهد چه نسبتی از نمونه‌های واقعی یک کلاس به درستی تشخیص داده شده‌اند:

\begin{equation}
Recall
=
\frac{TP}{TP+FN}
\end{equation}


% -------------------------------------------------
\subsection{\lr{F1-Score}}
% -------------------------------------------------

معیار
\lr{F1-Score}
میانگین هارمونیک
\lr{Precision}
و
\lr{Recall}
است:

\begin{equation}
F1
=
2
\frac{
Precision
\times
Recall
}{
Precision+Recall
}
\end{equation}


% -------------------------------------------------
\subsection{میانگین‌های \lr{Macro} و \lr{Weighted}}
% -------------------------------------------------

در میانگین
\lr{Macro}،
معیار برای هر کلاس به‌صورت مستقل محاسبه می‌شود و تمام کلاس‌ها وزن یکسانی دارند:

\begin{equation}
MacroPrecision
=
\frac{1}{K}
\sum_{k=1}^{K}Precision_k
\end{equation}

\begin{equation}
MacroRecall
=
\frac{1}{K}
\sum_{k=1}^{K}Recall_k
\end{equation}

\begin{equation}
MacroF1
=
\frac{1}{K}
\sum_{k=1}^{K}
F1_k
\end{equation}

در میانگین وزنی، سهم هر کلاس متناسب با تعداد نمونه‌های واقعی آن است. برای مثال:

\begin{equation}
WeightedF1
=
\sum_{k=1}^{K}
\frac{n_k}{N}F1_k
\end{equation}

که در آن $n_k$ تعداد نمونه‌های کلاس $k$ و $N$ تعداد کل نمونه‌ها است.


% -------------------------------------------------
\subsection{\lr{Confusion Matrix}}
% -------------------------------------------------

ماتریس درهم‌ریختگی\footnote{\lr{Confusion Matrix}} زبان واقعی نمونه‌ها را با زبان پیش‌بینی‌شده مقایسه می‌کند.

اگر $C_{ij}$ عضو سطر $i$ و ستون $j$ باشد:

\begin{equation}
C_{ij}
=
\text{تعداد نمونه‌های کلاس واقعی }i
\text{ که به کلاس }j
\text{ نسبت داده شده‌اند}
\end{equation}

مقادیر روی قطر اصلی نشان‌دهنده پیش‌بینی‌های صحیح و مقادیر خارج از قطر بیانگر خطاهای مدل هستند.


% =================================================
\section{معیارهای ارزیابی خوشه‌بندی}
% =================================================

معیارهای خوشه‌بندی به دو دسته داخلی و خارجی تقسیم می‌شوند. معیارهای داخلی فقط از ویژگی‌ها و تخصیص خوشه‌ها استفاده می‌کنند. معیارهای خارجی، مانند
\lr{ARI}،
\lr{NMI}
و
\lr{Purity}،
فقط پس از تشکیل خوشه‌ها و برای مقایسه با برچسب‌های واقعی به کار می‌روند.


% -------------------------------------------------
\subsection{\lr{Silhouette Score}}
% -------------------------------------------------

معیار
\lr{Silhouette Score}
برای بررسی انسجام داخل خوشه‌ها و میزان جدایی میان خوشه‌های مختلف استفاده می‌شود.

برای نمونه $i$:

\begin{itemize}
	\item $a(i)$ میانگین فاصله نمونه با اعضای خوشه خودش؛
	\item $b(i)$ کمترین میانگین فاصله نمونه با یک خوشه دیگر است.
\end{itemize}

در نتیجه:

\begin{equation}
s(i)
=
\frac{
b(i)-a(i)
}{
\max(a(i),b(i))
}
\end{equation}

و داریم:

\begin{equation}
-1
\leq
s(i)
\leq
1
\end{equation}

مقادیر نزدیک به 1 بیانگر کیفیت بهتر خوشه‌بندی هستند.

\subsection{\lr{Calinski-Harabasz Index}}

معیار
\lr{Calinski-Harabasz}
نسبت پراکندگی بین خوشه‌ها به پراکندگی درون خوشه‌ها را اندازه‌گیری می‌کند.

مقدار بالاتر این معیار نشان‌دهنده جدایی بهتر خوشه‌ها و فشردگی بیشتر نمونه‌های داخل هر خوشه است.

\subsection{\lr{Davies-Bouldin Index}}

معیار
\lr{Davies-Bouldin}
میزان شباهت میان هر خوشه و نزدیک‌ترین خوشه دیگر را بررسی می‌کند.

مقدار کمتر این معیار نشان‌دهنده کیفیت بهتر خوشه‌بندی است.
% -------------------------------------------------
\subsection{\lr{Cluster Purity}}
% -------------------------------------------------

معیار
\lr{Purity}
بررسی می‌کند اعضای هر خوشه تا چه اندازه به یک کلاس واقعی تعلق دارند.

اگر مجموعه خوشه‌ها برابر

\begin{equation}
\Omega
=
\{\omega_1,\ldots,\omega_K\}
\end{equation}

باشد،
\lr{Purity}
به صورت زیر تعریف می‌شود:

\begin{equation}
Purity(\Omega,C)
=
\frac{1}{N}
\sum_k
\max_j
|\omega_k\cap c_j|
\end{equation}

مقدار بیشتر
\lr{Purity}
نشان‌دهنده تطابق بیشتر خوشه‌ها با زبان‌های واقعی است.

این معیار به تعداد خوشه‌ها حساس است و ممکن است با تشکیل خوشه‌های کوچک‌تر افزایش یابد؛ بنابراین باید در کنار معیارهای دیگر تفسیر شود.


% -------------------------------------------------
\subsection{\lr{Adjusted Rand Index}}
% -------------------------------------------------

در بخش خوشه‌بندی، از
\lr{Adjusted Rand Index}
برای مقایسه ساختار خوشه‌های ایجادشده با برچسب‌های واقعی استفاده شده است.

فرم کلی این معیار به صورت زیر است:

\begin{equation}
ARI
=
\frac{
RI-E[RI]
}{
\max(RI)-E[RI]
}
\end{equation}

مقدار نزدیک به 1 بیانگر تطابق بالا میان خوشه‌بندی و برچسب‌های واقعی است.

مقدار یک نشان‌دهنده انطباق کامل، مقدار نزدیک به صفر نشان‌دهنده توافق مشابه حالت تصادفی و مقدار منفی نشان‌دهنده توافق کمتر از انتظار تصادفی است.


% -------------------------------------------------
\subsection{\lr{Normalized Mutual Information}}
% -------------------------------------------------

معیار
\lr{NMI}\footnote{\lr{Normalized Mutual Information}}
میزان اطلاعات مشترک میان تخصیص خوشه‌ها و برچسب‌های واقعی را اندازه‌گیری می‌کند. در نرمال‌سازی حسابی می‌توان نوشت:

\begin{equation}
NMI(C,Y)
=
\frac{2I(C;Y)}{H(C)+H(Y)}
\end{equation}

که در آن $I(C;Y)$ اطلاعات متقابل و $H(C)$ و $H(Y)$ آنتروپی خوشه‌ها و برچسب‌های واقعی هستند. مقدار این معیار در بازه صفر تا یک قرار می‌گیرد و مقدار بزرگ‌تر نشان‌دهنده انطباق بیشتر است.


% =================================================
\section{خلاصه و جمع‌بندی}
% =================================================

در این فصل مفاهیم نظری مورد نیاز برای درک مراحل پیاده‌سازی پروژه بررسی شدند. ابتدا مسئله شناسایی زبان از متن و گفتار و تفاوت میان یادگیری نظارت‌شده و بدون نظارت معرفی شد.

در بخش متن، روش‌های استخراج ویژگی شامل
\lr{Character N-Gram}،
\lr{Word N-Gram}،
\lr{TF-IDF}
و انتخاب ویژگی با
\lr{Chi-Square}
تشریح شدند. در بخش گفتار نیز نحوه استخراج ویژگی‌های
\lr{MFCC}،
\lr{Delta MFCC}،
\lr{Delta-Delta MFCC}،
ویژگی‌های طیفی، انرژی و مدت زمان بررسی شدند. تعداد و ساختار دقیق ویژگی‌های نسخه نهایی در فصل سوم ارائه خواهند شد.

همچنین الگوریتم‌های طبقه‌بندی و خوشه‌بندی مورد استفاده، روش‌های استانداردسازی و کاهش ابعاد،
\lr{Cross Validation}،
\lr{Grid Search}
و معیارهای ارزیابی معرفی شدند.

مفاهیم ارائه‌شده در این فصل، مبنای نظری مراحل آماده‌سازی داده، پیاده‌سازی مدل‌ها و تحلیل نتایج در فصل‌های بعدی خواهند بود.


